\section{Cash Equity and the Split Phantom}\label{sec:tour-equity}

The plainest product makes the plainest case, so the tour opens on cash equity: a
purchase, its cash dividend, and then the split that is the first event with a way to
go quietly wrong.

\subsection{The trade and its dividend}

The trade is where the ACME line begins. The book buys 1{,}000 ACME shares at
\EUR{100}, and the executed trade is a contract triggered by execution: it produces
one transaction, and that one transaction, at trade date, books both legs at once. It
brings 1{,}000 of the ACME share unit from the counterparty's virtual wallet into the
book's wallet, and it moves \EUR{100{,}000} of cash the other way, into the
counterparty's virtual wallet --- whose balance is now the unsettled payable the book
owes. Both legs are booked the instant the trade is struck, because on the trade date
the position is already owned and the cash is already due; what has not yet happened is
the external transfer that will hand the counterparty its money. The door checks that
ACME is a registered unit and that the transaction is well-formed, appends it, and the
book's share balance is 1{,}000 --- not a stored position some later process must
maintain, but the sum of the one move the log records --- while the counterparty's
virtual cash balance is the \EUR{100{,}000} still owed. Settlement, when it comes,
generates no moves at all: the positions were already correct from trade date, and the
external transfer of the cash is recorded as state --- the payable's confirmation,
marked on the transaction the log already holds --- not as a second transaction. The trade-date
booking and the settlement-date transfer are two dated facts joined by that
virtual-wallet balance, not a single instant in which shares and cash change hands
together; conservation holds within the one trade-date transaction, and nothing new is
needed to express it.

The cash dividend is the simplest lifecycle contract there is, and it is worth seeing
in full precisely because there is so little to it. ACME declares a \EUR{2} ordinary
dividend. This is the boundary $(f, c) = (1, 2)$: the quantity factor is 1, so no
share count changes, and the per-share cash is \EUR{2}. The contract, triggered by
the confirmed notice, reads the terms and the holding and submits one transaction. It
pays cash $1{,}000 \times 2 = \EUR{2{,}000}$ into the book's wallet, and it records
the declared operator \texttt{Shift} $-2$ as the boundary the read seam will
thereafter apply to prints: a cum-dividend quote of \EUR{102}, read across that
boundary, transports to $(102 - 2)/1 = \EUR{100}$, exactly the price going
ex-dividend. The stored price reference on the instrument itself is not moved --- under
an ordinary dividend a price terms field stands, and it is a market print crossing the
seam, not a stored field, that the boundary shifts. Value is conserved across the
boundary up to the cash the boundary paid:
\EUR{102{,}000} of shares beforehand is \EUR{100{,}000} of shares plus \EUR{2{,}000}
of cash after. There is no separate dividend engine. There is a contract that reads a
boundary and emits moves, and it is the same shape of contract that will, in a
moment, process a split.

A word on the feeds, because this basis convention is the ledger's own. Exchanges and
vendors rebase a quote series for a split or a large special dividend, but not for an
ordinary one: a print taken the morning after an ordinary ex-date arrives
\emph{unadjusted}, in the same basis as the day before. Treating the \EUR{2} ordinary
dividend as a boundary that carries a pre-ex quote to its ex-frame is therefore a
convention the ledger declares and applies at the read seam, reconciled there against
feeds that do not share it --- not a frame the feed already speaks. Nor does the
convention forecast the market: the realised ex-drop is nearer $c\,(1-\tau)$ than the
full $c$, because the dividend is taxed in the holder's hands, and that gap between the
transported quote and the next fresh print is ordinary marking profit and loss, never
something the boundary absorbs.

\subsection{The split, and the phantom it invites}

On the first of June, ACME splits two-for-one. The corporate-action contract reads
the confirmed notice, whose declaration is \texttt{Scale 2} --- the boundary
$(f, c) = (2, 0)$ --- and submits one transaction that does three things at once: it
moves the holding from 1{,}000 to 2{,}000 shares, it appends a new terms version
carrying the halved reference price of \EUR{50}, and it writes the new basis. Because
it is one transaction, no reader can ever see the shares doubled while the basis still
says pre-split, or the basis advanced while the shares have not moved. The door's
invariance weld admits it only if the quantity leg realises $q \mapsto q \cdot f$ with
no cash leg, which \texttt{Scale 2} with $c = 0$ does.

Now the event that breaks careless systems. At 09:59, one minute before the split
takes effect, a vendor prints an ACME price of \EUR{100}, stamped in the pre-split
basis. At 10:03 someone values the position. The share count is now 2{,}000. The most
recent print is \EUR{100}. Multiply them and the book is worth
$2{,}000 \times 100 = \EUR{200{,}000}$. Every factor here is individually correct:
2{,}000 is the true share count, and \EUR{100} was a true, recently printed price.
The product of two correct numbers, \EUR{200{,}000}, contains \EUR{100{,}000} of
phantom value on top of the true \EUR{100{,}000} --- and no internal reconciliation can
catch it, because nothing internal disagrees: the position system and the valuation
system are both reading real data. The error is that the two
numbers live in different frames: the quantity is post-split, the price is pre-split,
and they were never entitled to meet.

The right number is \EUR{100{,}000}, and it can be reached from either frame. In the
pre-split frame, $1{,}000 \times 100 = \EUR{100{,}000}$. In the post-split frame, the
stale print is first transported across the boundary by the same declared
\texttt{Scale 2}, now consumed at the read seam: $100 / 2 = \EUR{50}$, and only then
allowed to meet the 2{,}000-share count, $2{,}000 \times 50 = \EUR{100{,}000}$. Both
frames agree because there is one declared operator and it is applied wherever a
pre-split number tries to cross. Single-basis consumption is the rule the door
enforces here: a quantity and a price may be multiplied only in one frame, and it is
the ledger that decides the frame, from the declared boundary, not the feed that
happens to have printed last.

The design forbids the phantom by refusing to build it. The 09:59 print is not
deleted or corrected --- it is real data and stays in the log --- but the pairing of a
pre-split price with a post-split quantity is not a wrong number the door computes and
lets through. It is a transaction that does not typecheck, and a stale datum met at a
boundary with no declared conversion is a typed refusal, not a silent \EUR{200{,}000}.

One more turn shows why the frame has to be decided by declared order and not by
arrival order. Suppose the book pulls a history a week later, and reads a May print of
\EUR{102} across both the May dividend and the June split. The two boundaries compose
in effective order: dividend first, then split, giving $(102 - 2)/2 = \EUR{50}$. The
reverse order, $102/2 - 2 = \EUR{49}$, is a different number, and it is not a valid
conversion of anything: no pair of bases on the chain has it, and no operator the
menu can build will forge it. The arithmetic does not commute, so if arrival order
could pick the order of composition, the same datum would convert to \EUR{50} or
\EUR{49} depending on when it happened to be read. It cannot, because the composition
is fixed by the declared boundaries, and \EUR{50} is the only number the declared
data admits.

None of this required restating a single stored share. The 1{,}000 became 2{,}000 by
one recorded move, and the reference price became \EUR{50} by one appended terms
version; the old count and the old price both still stand in the log exactly as
written. \emph{The decision this stop rests on is single-basis consumption composed in
declared order: a quantity and a price meet only in one frame, the ledger fixes that
frame from the declared boundary rather than from whichever feed printed last, and when
boundaries stack it is the declared effective order --- not the order a datum happens to
arrive --- that fixes the single number it converts to, which is why \EUR{50} is forced
and \EUR{49} cannot be assembled. That quantities are never restated in place is the
visible consequence, not the cause: the phantom is unrepresentable because the pairing
that would build it does not typecheck, not because a restatement was performed
carefully.}
